\chapter{Experimentos}
\label{chapter:experimentos}

Para validar o funcionamento da plataforma contra um ambiente Proxmox VE e um cluster k3s reais — e não apenas contra simulações —, foram conduzidas validações de fluxo de trabalho, executadas tanto por utilitários de linha de comando quanto pelo uso ponta a ponta do portal web. Cada validação percorre um fluxo completo, do ponto de vista do usuário, e verifica se o resultado esperado é observado no ambiente real.

\section{Validação de criação de recursos}

Foram implementados dois utilitários de linha de comando (\texttt{cmd/proxmox-check}) que exercitam o fluxo completo de provisionamento contra a infraestrutura real:

\begin{itemize}
    \item o primeiro cria uma máquina virtual por clone do \textit{template} Ubuntu com \textit{cloud-init} (o mesmo caminho utilizado pelo portal), consulta o seu status e a remove, validando a integração via Terraform e via API REST;
    \item o segundo, executado com a opção \texttt{--k3s}, cria, consulta e remove um contêiner de aplicação no cluster k3s real, validando a integração via \texttt{client-go}.
\end{itemize}

\par Em seguida, o mesmo fluxo foi executado pelo portal web, com a criação de uma máquina virtual Ubuntu — acessada via terminal web e removida — e de um contêiner de aplicação, conforme a Seção~\ref{sec:val_conteudo}. Foram verificados também os limites por perfil: a criação de VM por um usuário com perfil \textit{aluno} é feita no \textit{pool} definido pelo \textit{backend} para esse perfil, independentemente do valor enviado pelo cliente.

\figimg{fig:val_vm_proxmox}{Máquina virtual criada pela plataforma, vista na interface do Proxmox VE, com o histórico de tarefas (\textit{resize}, \textit{start} e a remoção da VM do teste de fumaça)}{val_vm_proxmox}{captura de tela do Proxmox VE}

\figimg{fig:val_saida_smoke}{Saída dos testes de fumaça contra o ambiente real: contêiner no cluster k3s (\texttt{--k3s}) e máquina virtual no Proxmox VE}{val_saida_smoke}{captura de tela do terminal}

\section{Validação de criação de conteúdo}
\label{sec:val_conteudo}

% TODO(autor): confirmar se "criação de conteúdo" no sentido do professor é o que foi feito dentro dos recursos (banco, dados, logs). Ajustar o texto se for outro entendimento.
Além de criar o recurso, validou-se que é possível gerar e utilizar conteúdo dentro dele. Foi criado, pelo portal, um contêiner de aplicação PostgreSQL com variáveis de ambiente, volume persistente e exposição de porta TCP via \textit{NodePort}; foram consultados os logs do contêiner e validada uma conexão via \texttt{psql}, na qual o banco respondeu às consultas, antes de o contêiner ser removido. De modo análogo, na VM Ubuntu, o terminal web foi utilizado para executar comandos no sistema criado.

\par O acesso externo a um contêiner por meio do par \texttt{ip:porta} do \textit{NodePort} não pôde ser validado durante o desenvolvimento, pois a rede VPN utilizada para acesso ao ambiente de laboratório bloqueia o tráfego de entrada nessas portas a partir de fora da rede interna. Essa validação depende de o \textit{backend} estar em execução na própria VM do laboratório, e é discutida como trabalho futuro no Capítulo~\ref{chapter:conclusoes_trabalhos_futuros}.

\figimg{fig:val_postgres}{Detalhes do contêiner PostgreSQL em execução, com endereço externo (\textit{NodePort}) e variáveis de ambiente}{front-tela-container-em-execucao}{captura de tela do ProxLab}

\figimg{fig:val_logs_container}{Terminal e logs do contêiner PostgreSQL: conexão via \texttt{psql} e listagem das tabelas criadas}{front-tela-terminal-container}{captura de tela do ProxLab}

\section{Validação de \textit{pass-through} de GPU}

% TODO(autor): GPU NÃO está implementada no código (fora da v1). Se o professor quer este experimento, é preciso
% implementar/testar o pass-through (PCIe passthrough no Proxmox / hostpci no Terraform) antes de preencher esta seção.
% Se não for possível até 15/11, manter como trabalho futuro e remover esta seção.
Esta validação tem por objetivo verificar se uma máquina virtual criada pela plataforma consegue acessar uma GPU física do servidor por meio de \textit{pass-through} de dispositivo PCIe, requisito de cargas de IA/ML. Até a presente versão, o acesso a GPU não foi implementado, razão pela qual não há resultados a reportar; o tema é tratado como trabalho futuro no Capítulo~\ref{chapter:conclusoes_trabalhos_futuros}.

\figplaceholder{fig:val_gpu}{GPU visível dentro da máquina virtual (\texttt{nvidia-smi})}{Captura do terminal da VM exibindo a saída de \texttt{nvidia-smi} (a ser feita somente se o pass-through for implementado)}{captura de tela do terminal}

\section{Validação de criação e autenticação de usuários}

Validou-se o fluxo de identidade da plataforma: (i) o \textit{login} retorna um \textit{token} JWT e as rotas protegidas rejeitam requisições sem \textit{token} ou com \textit{token} inválido (resposta 401, com redirecionamento para a tela de \textit{login}); (ii) o cadastro de um novo usuário só é aceito quando solicitado por um usuário com perfil \textit{pesquisador}, sendo recusado para o perfil \textit{aluno}; (iii) o usuário consegue alterar seu nome e sua senha, sendo recusada a troca quando a senha atual informada está incorreta; e (iv) o pesquisador consegue listar, alterar o e-mail e excluir usuários, sendo a exclusão acompanhada da remoção das VMs e dos contêineres do usuário.

\figimg{fig:val_usuarios}{Tela de cadastro de novo usuário pelo pesquisador}{front-tela-cadastro-usuario}{captura de tela do ProxLab}
